本文构建包含56,000条LaTeX公式的数学分词benchmark，按领域、来源和复杂度分类。我们研究数学原子BPE与软AST-BPE，并与标准AST-BPE及五个基准比较。实验包含领域内、跨领域、极限泛化和简单到复杂四类设置，使用六项指标评价压缩、数学单位完整性、语法跨越、词表使用与编码速度。软AST-BPE在11组划分中均取得最低T/B，领域内相对BPE减少7.45\%的每字节token数，命令完整率为99.73\%。标准AST-BPE的Cross更低，软AST-BPE则获得更好的压缩与编码吞吐。
